\documentclass[10pt, a4paper, twocolumn]{article}

\usepackage[utf8]{inputenc}
\usepackage[T1]{fontenc}
\usepackage{geometry}
\usepackage{hyperref}
\usepackage{amsmath}
\usepackage{booktabs}
\usepackage{microtype}
\usepackage{parskip}

\geometry{margin=1.8cm, top=2cm}

\title{\textbf{Image Fragment Reconstruction via Adjacency Prediction}\\[4pt]
\large Case Study Report}
\author{}
\date{}

\begin{document}

\maketitle
\thispagestyle{empty}

% ─────────────────────────────────────────────────────────────────────────────
\section*{1. Model Architecture}
% ─────────────────────────────────────────────────────────────────────────────

The model is a \textbf{Siamese CNN} that operates on pairs of $16{\times}16$
fragments. A shared encoder maps each fragment to a 256-dimensional embedding;
a single comparison head maps the two embeddings to a scalar confidence score
$p_{ij} \in [0,1]$.

\paragraph{Encoder.}
Three Conv--BatchNorm--ReLU blocks with max-pooling after the last two reduce
the spatial dimensions from $16{\times}16$ to $4{\times}4$; a linear layer
projects the flattened features to 256-d (ReLU + dropout $p=0.3$).

\paragraph{Comparison head.}
Inspired by InferSent~\cite{infersent2017}, the head takes the
1024-dimensional vector
$[\mathbf{a}-\mathbf{b}\;\|\;\mathbf{a}\odot\mathbf{b}\;\|\;\mathbf{a}\;\|\;\mathbf{b}]$
as input, passes it through one hidden layer (ReLU), and produces a single
scalar via sigmoid. Using the difference and Hadamard product alongside the
raw embeddings gives the head explicit access to asymmetric relational
features at no additional encoder cost.

\paragraph{Why Siamese?}
Weight sharing ensures the comparison is symmetric and forces the encoder to
learn a universal fragment representation rather than one tied to a specific
input slot. The pairwise architecture avoids the need for any image-level
label: the only supervision required is whether two fragments are adjacent,
which is derivable from the grid structure alone.

% ─────────────────────────────────────────────────────────────────────────────
\section*{2. Training Strategy}
% ─────────────────────────────────────────────────────────────────────────────

\paragraph{Pretext task.}
The model is trained to predict spatial adjacency: $y_{ij}=1$ iff fragments
$i$ and $j$ are grid-neighbours within the same image. With $N=10$ images on
a $4{\times}4$ grid, the negative-to-positive ratio is exactly
$12{,}480/240 = 52$. We use weighted BCE with tilt parameter $\beta$:
\[
  \mathcal{L}_{\text{WBCE}}(\beta) = -\tfrac{1}{N}\textstyle\sum_{(i,j)}
  \bigl[\beta\cdot 52\, y_{ij}\log p_{ij} + (1-y_{ij})\log(1-p_{ij})\bigr].
\]

\paragraph{Multi-task loss.}
Adjacency is a local signal; the downstream task is global clustering.
A same-image loss term $\mathcal{L}_{\text{BCE}}^{\text{same}}$ is added to
bridge this gap, acting on the same score $p_{ij}$:
\[
  \mathcal{L} = \lambda_{\text{adj}}\,\mathcal{L}_{\text{WBCE}}(\beta)
              + \lambda_{\text{same}}\,\mathcal{L}_{\text{BCE}}^{\text{same}}.
\]

\paragraph{Inference.}
The $160{\times}160$ matrix of scores $p_{ij}$ is used directly as a
similarity matrix for \textbf{balanced spectral clustering}: spectral
embedding (normalised graph Laplacian, $k=10$ components) followed by
Hungarian-assignment balanced $k$-means that enforces exactly 16 fragments
per cluster.

\paragraph{Training details.}
1.28M training images (ImageNet-64); each step samples 10 fresh images.
Up to 25{,}000 iterations; early stopping with patience 10, evaluated every
250 steps over 20 validation batches. All runs on CPU (Apple M-series).

% ─────────────────────────────────────────────────────────────────────────────
\section*{3. Evaluation Results}
% ─────────────────────────────────────────────────────────────────────────────

Performance is measured by the \textbf{Adjusted Rand Index} (ARI), which
counts pair-level agreement between predicted and true clusters, adjusted for
chance. ARI\,=\,0 is random; ARI\,=\,1 is perfect.

\paragraph{Validation results.}
A five-seed study (seeds $0$--$4$) across five hyperparameter configurations
(Table~\ref{tab:results}) shows ARI in the range $0.50$--$0.58$, stable
across seeds ($\sigma\approx 0.02$--$0.03$). A Welch one-way ANOVA finds no
statistically significant difference between configurations
($\alpha=0.10$), indicating the choice of loss weights is not a strong lever
at this scale.

\begin{table}[h]
\centering\small
\begin{tabular}{lcc}
\toprule
\textbf{Configuration} & \textbf{Mean ARI} & \textbf{SD} \\
\midrule
(a) adj-only, $\beta=1/52$     & TODO & TODO \\
(b) same-only                  & TODO & TODO \\
(c) multi-task $\lambda_s=1.0$ & TODO & TODO \\
(d) multi-task $\lambda_s=0.2$ & TODO & TODO \\
(e) adj-only, $\beta=1$        & TODO & TODO \\
\bottomrule
\end{tabular}
\caption{Validation ARI (5 seeds each). % TODO fill after overnight runs.
}
\label{tab:results}
\end{table}

\paragraph{Test results.}
% TODO: fill after running bash scripts/run_test_eval.sh <best_run>
The best checkpoint (run \texttt{TODO}) was evaluated on 1{,}000 independent
test batches:
ARI\,=\,TODO\,$\pm$\,TODO,
NMI\,=\,TODO\,$\pm$\,TODO,
purity\,=\,TODO\,$\pm$\,TODO,
adjacency F1\,=\,TODO.

\paragraph{Pretext task.}
Adjacency AUROC sits at $0.95$--$0.97$ across all configurations, indicating
the model has nearly solved the binary pretext task. The remaining gap to
perfect clustering is attributed to fragments with near-uniform colour or
repetitive texture that look similar across different source images.

% ─────────────────────────────────────────────────────────────────────────────
\section*{4. Challenges and Improvements}
% ─────────────────────────────────────────────────────────────────────────────

\paragraph{Class imbalance.}
The 52:1 negative-to-positive ratio in adjacency pairs requires careful loss
weighting. We found $\beta=1/52$ (plain BCE) works as well as the textbook
class-balanced $\beta=1$; heavier upweighting degrades performance.

\paragraph{Local vs.\ global signal.}
Adjacency is a nearest-neighbour signal; two fragments from opposite corners
of the same image are labelled negative. The same-image loss term partially
closes this gap but does not produce a statistically significant improvement
at this scale, suggesting the binding constraint is fragment ambiguity rather
than the loss formulation.

\paragraph{Fragment ambiguity.}
Uniform-colour or repetitive-texture patches yield near-identical embeddings
regardless of source, fundamentally limiting the achievable ARI for any
pairwise model without global context.

% ─────────────────────────────────────────────────────────────────────────────
\section*{5. Next Steps}
% ─────────────────────────────────────────────────────────────────────────────

\paragraph{Pretrained encoder.}
Replacing the trained-from-scratch CNN with a self-supervised encoder
(e.g.\ DINOv2) would provide richer features at $16{\times}16$ resolution
at no labelling cost, likely the highest-impact single change.

\paragraph{Topology-aware clustering.}
The true clustering is more constrained than ``equal sizes'': each cluster's
induced subgraph must be isomorphic to a $4{\times}4$ grid. A GNN-based
approach that exploits this planarity constraint could resolve borderline
assignments that spectral clustering misses.

\paragraph{Larger-scale statistical study.}
With $n=5$ seeds the ANOVA is underpowered for effect sizes below
$\sim\!0.07$ ARI. A 20-seed study would give $\Delta_{\min}\approx0.03$,
sufficient to detect the differences suggested by the single-seed sweep.

\paragraph{Color-histogram baseline.}
A non-learned baseline using per-fragment color histograms with spectral
clustering would anchor the model's ARI against a trivial reference and
quantify how much benefit comes from learned representation versus color
matching.

\begin{thebibliography}{9}\small
\bibitem{infersent2017}
Conneau et al.\ (2017). \textit{Supervised Learning of Universal Sentence
Representations from Natural Language Inference Data}. EMNLP.
\end{thebibliography}

\end{document}
